\SourceLeaf{130}
It is easy to see how any number $A$ can be brought into the above form. The values of $\mu$ and $a$ are immediately apparent. If one then sets $A = 10^{\mu} a B$, $b$ is the digit in the first decimal place of $B$. If one further sets $B = \left(1+\frac{1}{10} b\right) C$, $c$ is the digit in the second decimal place of $C$, and so on. A closer look shows that it is still more convenient to bring $\frac{1}{A}$ into the above form without actually calculating this quotient, by setting successively $10^{\mu} a A = B$, $\left(1+\frac{1}{10} b\right) B = C$, and so on. Thus one obtains $B$ from $A$, $C$ from $B$, and so on, by multiplication, whereas the first method requires division. \textsc{Leonelli} has, however, further shortened the calculation in both methods by reducing it to a fixed procedure. Anyone expecting to use these tables frequently would do well to become familiar with it. After some practice, one will surely soon come to exercise the requisite attention automatically. We therefore agree with \textsc{Leonelli} that \textsc{Delambre}'s objection is hasty when he says that wearisome attention is necessary unless one writes down all those superfluous zeros every time.

It is easy to see that one can similarly bring $A$ or $\frac{1}{A}$ into the form
\[
10^{\mu} a \left(1+\frac{1}{100} b\right)\left(1+\frac{1}{10000} c\right)\ldots
\]
so that $a$, $b$, $c\ldots$ all denote integers below 100. In each individual case one will then need only half as many terms. The auxiliary table likewise needs only half as many orders, but 99 logarithms in each, and must therefore be almost five times as extensive. \textsc{Leonelli} has appended such an auxiliary table too, though only for Briggsian logarithms and only to 15 decimal places. It contains 485 logarithms. It is a pity that this table too has not been calculated to 20 decimal places.

\textsc{Leonelli}'s tables are no less applicable to the inverse problem of finding the number for a given logarithm. Breaking the given logarithm down into its constituents found in the tables gives the sequence of factors of the number. Multiplying them, for which a convenient algorithm is easily found, gives the number itself.

